\scp{Regional interlanguages}{15-05}
\citet{Schapper2020a} discusses the concentrated distribution
of reflexes of the pseudo-etymon {\#}lutuR ‘wall built up of
stones’ from the eastern tip of Timor across to Aru,
to a lesser degree in Flores,
but not much in central Maluku
and “not in evidence across the whole proposed CMP area,
[{\ldots}] limited to the southern arc of islands” \citep[241]{Schapper2020a}.
Reflexes are also found in Papuan Timor-Alor-Pantar languages.
Those are assumed to be from an Austronesian source.
The presence of doublets in single languages,
one of which displays regular sound correspondences,
and the other of which does not,
is seen as evidence of diffusion, rather than retention.
For example, Ujir \it{luti} ‘stone wall (n); stack up stones
(v)’ shows regular sound correspondences (inherited),
whereas Ujir \it{lutur} ‘stone wall fortification’ is irregular
and marks it as borrowed (diffusion) (p.242).
Similarly Central Marsela \it{lukra} ‘stone wall around garden’
shows regular sound correspondences (inherited),
whereas Central Marsela \it{lutur} ‘protective stone wall around
village’ is irregular and marks it as borrowed (diffusion) (p.242).
This uneven distribution through diffusion in some subgroups
but not others in Wallacea is similar to the other items discussed
in the previous sections.

\citet{Borges2014} provides a parallel scenario of contact between
two completely unrelated groups that may help explain the distribution
of words in Wallacea which are shared by some subgroups but not others.
The indigenous languages of the South American country of Suriname are Amerindian.
There are other, completely unrelated languages known loosely
as ``Maroon'', that developed among groups of African slaves,
many of whom escaped into the interior immediately after their
arrival and preserved many aspects of their native (mostly west) African languages and cultures.
\citet[99]{Borges2014} summarises the situation as follows:

\begin{quote}
	“The earliest Surinamese Maroons,
	for which we have records,
	appear not to have promulgated any Creole language,
	rather they assimilated culturally
	and linguistically to the Amerindian groups they encountered.
	Other, slightly later, early Maroons intermarried with Amerindian
	women, thereby adopting established material culture
	and techniques for surviving in the local environment.
	As they took over things
	and tactics for which they had no words,
	the Maroons borrowed Amerindian words
	and incorporated them into their incipient languages.
	The development of an intense trade relationship between the
	Ndyuka and their upriver Amerindian neighbors in the mid-18th century,
	combined with the lack of a functional interlanguage,
	led to the most extreme linguistic effect -- the development of Ndyuka-Trio Pidgin.
	From the 1960s, however,
	the pidgin has ceased to be used due to a decline in the intensity
	of direct trade between the groups.” \citep[99]{Borges2014}
\end{quote}

To recap and paraphrase, there are three linguistic scenarios that developed.
Keep in mind that the two unrelated groups each have significant internal variety.
\it{Scenario {\#}1}: some newcomers (the Maroons [read: Austronesians])
assimilated culturally and linguistically to become like the established traditional
owners (the Amerindians [read: Papuans]).
\it{Scenario {\#}2}: some newcomers intermarried with the traditional
owners, and adopted into their own languages new words
and techniques appropriate to the new environment
and exposure to other ways of doing
and talking about things.
\it{Scenario {\#}3}: the necessity of intense trade
and the need to communicate eventually led to the development
of an interlanguage (shared language of wider communication, trade language).
And of particular interest to us in Wallacea is that \citet{HuttarVelantie1997}
pointed out that Ndyuka-Trio Pidgin ceased to be used by around
1960, when contact diminished and the utility of other interlanguages
superseded the need for its continued use.

The time-depth in Suriname is relatively short,
focusing on a period of contact of only about 300 years (1660--1960),
with the useful role of Ndyuka-Trio Pidgin narrowed down to 200 years or so.
By contrast, the time-depth of contact in Wallacea between Austronesian-speaking
peoples and Papuan-speaking peoples spans about 3,500--4,000 years,
more than enough time for interlanguages to both develop
and also to disappear.

In previous chapters we have noted a number of languages that
have functioned as interlanguages,
trade languages or languages of wider communication in a limited regional role.
For example, we noted that the STB language of Geser-Gorom has
functioned for centuries as a language of wider communication
in the regions of eastern Seram
and the Bomberai peninsula (a.k.a.
Bird’s Neck) of west Papua
\citep{LoskiLoski1989,WalkerWalker1991}.
It was used in the trade links with Onin,
traces of Geser loans are widespread in the region in both Austronesian
and Papuan languages, and its role has waned significantly in
the last 60 years or so as its utility has been superceded by Malay.

Similarly, the STB language of Fordata has functioned in the
past as a language of wider communication
and the favoured language of ceremony in the Kei islands,
and in Tanimbar in the Yamdena-speaking regions \citep{Hughes1987,McKinnon1991}.
Traces are still found in Yamdena rituals,
but its role has waned significantly.

We noted that the \tsc{North Halmahera} (West Papuan) language
of Ternate has left significant traces throughout the Sula archipelago,
Buru, western Seram, and in languages spoken in the Manipa Straits
between Buru and Seram \citep{Collins1982b,Grimes1993}.
Ambon Malay \it{kusu} ‘cuscus’ is very likely a Ternate loan
from \it{kuso} ‘cuscus’ (regular from Proto North Halmahera *kusor
‘cuscus’), with no similar forms being found in any other language of central
Maluku (see appendix \ref{ch:SurMar}).
During the height of the European spice trade the role of the
sultan of Ternate had significant reach into central Maluku,
the legacy of which is still felt today.
The \tsc{North Halmahera} (West Papuan) language of Tidore had
a similar influence in the Bird’s Head region of west Papua.

The \tsc{Southwest Maluku} (\tsc{Timor-Babar}) language of Luang
has functioned as a ceremonial language
and in other roles extending its reach both westward,
as well as to the Babar islands
and even the Kei islands far to the east \citep[400, note 13]{Taber1993}.

In \crf{ch:CT}, we noted that Tetun [tet] has had significant
impact on \tsc{Central Timor} languages over a long period of time.
More recently with the independence of Timor-Leste,
the different variety of Tetun Dili [tdt] has seen a resurgence
of being an important loan source into the Papuan
and Austronesian languages of central and eastern Timor.

We noted that prior to 1727,
Makasar (in South Sulawesi) had a considerable impact on Sumbawa
and western Flores, and after 1727 Bima exerted its influence
into western Flores \citep{VerheijenGrimes1995}.

There are widespread traces throughout languages in the region
of the waxing and waning in different roles for Javanese,
Makasar, Sanskrit, Arabic, Portuguese,
and Dutch, most of which spanned several centuries.
The role of Malay has gradually increased in the region over
the past thousand years or so,
and has superseded the role or need for many more localised languages of wider communication.
This is further complicated by the earlier role of trade Malay
that developed into a number of different contact creoles in
the region (such as Ambon Malay,
Kupang Malay, Aru Malay, Larantuka Malay, etc.).
And these are now in contact with Indonesian (Standard Malay),
which developed from official language/literary Malays of the
sultans’ courts in the western regions,
rather different from the trade Malay formerly used in what is
now eastern Indonesia \citep{AdelaarPrentice1996,Grimes1991,Grimes1996,Prentice1978}.

The languages mentioned above are only a sample,
but should be sufficient to make the point that there have been
many different languages of wider communication that have left
identifiable traces in the region,
and across different subgroups.
The languages mentioned above only account for the past 500 years or so.
But the time depth of around 4,000 years of contact should mean
that it is not unreasonable for there to have been other languages
of wider communication in more localised parts of the region
that we cannot identify,
and which no longer exist or no longer function in that wider role.
The distribution of some words across several subgroups may reflect
traces of such localised languages of wider communication.

This provides a possible explanation for why certain words are
found unevenly across several subgroups or nodes, but not others.
The CEMP hypothesis and diffusion does not explain the problems
with the uneven distribution of these lexical items.
Appealing to regional languages of wider communication,
for which we have some direct evidence
and some trace evidence,
and some of which may no longer exist,
is one reasonable way to account for the uneven distribution
of some lexical items across seemingly otherwise unrelated subgroups.

\newpage
While this does account for the complexities of the data we find
in the region, and does account for the problem that the distribution of one
word does not align with the distribution of other words,
it also flies in the face of the principles
and assumptions behind the classical Comparative Method -- a method
which is known to break down in contact situations.
Precisely the situation we find in eastern Indonesia and Timor-Leste.
And not only is there contact with languages of other families
and typologies, but also extensive trade spanning millennia
and a high degree of mobility throughout the region (see \srf{02-02}).

The Blust line loosely marks \it{a threshold into a transition
zone of contact} -- in which a number of features are unevenly
distributed, but collectively tell a similar story: around 4000--3,500 years
BP, Austronesian-speaking peoples moved into
and through the region fairly rapidly; the Austronesian languages
they spoke diverged very little from PMP in their phonology
and lexicon; there were probably several incursions from several
places and possibly in different time periods; they did not advance
in a single wave; they got off their boats in different places,
sometimes leap-frogging past islands where other groups stopped
off; they encountered pre-existing populations in place speaking
a variety of typologically SOV Papuan languages; the Austronesian-speaking
peoples were accommodating and receptive, learned the Papuan languages,
became bilingual and intermarried; words were borrowed for talking
about new and strange things, such as marsupials
and certain food crops; over time the Austronesian languages
adapted in some parts of the grammar
and lexicon to function more like the Papuan languages.
And conversely, many speakers of Papuan languages learned Austronesian
languages, there was intermarriage, and some Papuan languages adapted to
function more like some features of Austronesian languages.
Several regional interlanguages came and went.
There were varying degrees of language mixing (see \citealp{Ross2017}), in both directions.